\PassOptionsToPackage{unicode}{hyperref}
\PassOptionsToPackage{hyphens}{url}
\PassOptionsToPackage{dvipsnames,svgnames,x11names}{xcolor}
\documentclass[10pt,letterpaper]{article}

\usepackage{xcolor}
\usepackage[margin=2.0cm]{geometry}
\usepackage{lmodern}
\usepackage{amsmath,amssymb,amsfonts,mathtools}
\usepackage{textcomp}
\usepackage{microtype}
\usepackage{parskip}
\usepackage{longtable,booktabs,array,calc}
\usepackage{caption}
\captionsetup[table]{skip=6pt,font=small,labelfont=bf}
\captionsetup[figure]{font=small,labelfont=bf}
\newcounter{none}
\usepackage{etoolbox}
\usepackage{graphicx}
\usepackage{fancyvrb}
\newcommand{\VerbBar}{|}
\newcommand{\VERB}{\Verb[commandchars=\\\{\}]}
\DefineVerbatimEnvironment{Highlighting}{Verbatim}{commandchars=\\\{\},fontsize=\footnotesize}
\RecustomVerbatimEnvironment{verbatim}{Verbatim}{fontsize=\footnotesize}
\newenvironment{Shaded}{}{}
\newcommand{\AlertTok}[1]{\textcolor[rgb]{0.80,0.00,0.00}{\textbf{#1}}}
\newcommand{\AnnotationTok}[1]{\textcolor[rgb]{0.38,0.63,0.69}{\textbf{\textit{#1}}}}
\newcommand{\AttributeTok}[1]{\textcolor[rgb]{0.49,0.56,0.16}{#1}}
\newcommand{\BaseNTok}[1]{\textcolor[rgb]{0.25,0.63,0.44}{#1}}
\newcommand{\BuiltInTok}[1]{\textcolor[rgb]{0.00,0.50,0.00}{#1}}
\newcommand{\CharTok}[1]{\textcolor[rgb]{0.25,0.44,0.63}{#1}}
\newcommand{\CommentTok}[1]{\textcolor[rgb]{0.38,0.63,0.69}{\textit{#1}}}
\newcommand{\CommentVarTok}[1]{\textcolor[rgb]{0.38,0.63,0.69}{\textbf{\textit{#1}}}}
\newcommand{\ConstantTok}[1]{\textcolor[rgb]{0.53,0.00,0.00}{#1}}
\newcommand{\ControlFlowTok}[1]{\textcolor[rgb]{0.00,0.44,0.13}{\textbf{#1}}}
\newcommand{\DataTypeTok}[1]{\textcolor[rgb]{0.56,0.13,0.00}{#1}}
\newcommand{\DecValTok}[1]{\textcolor[rgb]{0.25,0.63,0.44}{#1}}
\newcommand{\DocumentationTok}[1]{\textcolor[rgb]{0.73,0.13,0.13}{\textit{#1}}}
\newcommand{\ErrorTok}[1]{\textcolor[rgb]{0.80,0.00,0.00}{\textbf{#1}}}
\newcommand{\ExtensionTok}[1]{#1}
\newcommand{\FloatTok}[1]{\textcolor[rgb]{0.25,0.63,0.44}{#1}}
\newcommand{\FunctionTok}[1]{\textcolor[rgb]{0.02,0.16,0.49}{#1}}
\newcommand{\ImportTok}[1]{\textcolor[rgb]{0.00,0.50,0.00}{\textbf{#1}}}
\newcommand{\InformationTok}[1]{\textcolor[rgb]{0.38,0.63,0.69}{\textbf{\textit{#1}}}}
\newcommand{\KeywordTok}[1]{\textcolor[rgb]{0.00,0.44,0.13}{\textbf{#1}}}
\newcommand{\NormalTok}[1]{#1}
\newcommand{\OperatorTok}[1]{\textcolor[rgb]{0.40,0.40,0.40}{#1}}
\newcommand{\OtherTok}[1]{\textcolor[rgb]{0.00,0.44,0.13}{#1}}
\newcommand{\PreprocessorTok}[1]{\textcolor[rgb]{0.74,0.48,0.00}{#1}}
\newcommand{\RegionMarkerTok}[1]{#1}
\newcommand{\SpecialCharTok}[1]{\textcolor[rgb]{0.25,0.44,0.63}{#1}}
\newcommand{\SpecialStringTok}[1]{\textcolor[rgb]{0.73,0.40,0.53}{#1}}
\newcommand{\StringTok}[1]{\textcolor[rgb]{0.25,0.44,0.63}{#1}}
\newcommand{\VariableTok}[1]{\textcolor[rgb]{0.10,0.09,0.49}{#1}}
\newcommand{\VerbatimStringTok}[1]{\textcolor[rgb]{0.25,0.44,0.63}{#1}}
\newcommand{\WarningTok}[1]{\textcolor[rgb]{0.38,0.63,0.69}{\textbf{\textit{#1}}}}

\usepackage{newunicodechar}
\newunicodechar{µ}{\ifmmode\mu\else\textmu\fi}
\newunicodechar{°}{\ifmmode^\circ\else\textdegree\fi}
\newunicodechar{·}{\ifmmode\cdot\else\textperiodcentered\fi}
\newunicodechar{²}{\ensuremath{^{2}}}
\newunicodechar{³}{\ensuremath{^{3}}}
\newunicodechar{¹}{\ensuremath{^{1}}}
\newunicodechar{⁴}{\ensuremath{^{4}}}
\newunicodechar{⁶}{\ensuremath{^{6}}}
\newunicodechar{⁻}{\ensuremath{^{-}}}
\newunicodechar{₀}{\ensuremath{_{0}}}
\newunicodechar{ₙ}{\ensuremath{_{n}}}
\newunicodechar{Δ}{\ensuremath{\Delta}}
\newunicodechar{Ω}{\ensuremath{\Omega}}
\newunicodechar{α}{\ensuremath{\alpha}}
\newunicodechar{β}{\ensuremath{\beta}}
\newunicodechar{δ}{\ensuremath{\delta}}
\newunicodechar{ε}{\ensuremath{\varepsilon}}
\newunicodechar{θ}{\ensuremath{\theta}}
\newunicodechar{κ}{\ensuremath{\kappa}}
\newunicodechar{λ}{\ensuremath{\lambda}}
\newunicodechar{π}{\ensuremath{\pi}}
\newunicodechar{ρ}{\ensuremath{\rho}}
\newunicodechar{σ}{\ensuremath{\sigma}}
\newunicodechar{τ}{\ensuremath{\tau}}
\newunicodechar{→}{\ensuremath{\rightarrow}}
\newunicodechar{∎}{\ensuremath{\blacksquare}}
\newunicodechar{−}{\ensuremath{-}}
\newunicodechar{∧}{\ensuremath{\wedge}}
\newunicodechar{∫}{\ensuremath{\int}}
\newunicodechar{≈}{\ensuremath{\approx}}
\newunicodechar{≤}{\ensuremath{\leq}}
\newunicodechar{≥}{\ensuremath{\geq}}
\newunicodechar{⊆}{\ensuremath{\subseteq}}
\newunicodechar{✓}{\ensuremath{\checkmark}}
\newunicodechar{≠}{\ensuremath{\neq}}
\newunicodechar{↔}{\ensuremath{\leftrightarrow}}
\newunicodechar{≪}{\ensuremath{\ll}}
\newunicodechar{↛}{\ensuremath{\nrightarrow}}
\newunicodechar{⇒}{\ensuremath{\Rightarrow}}
\newunicodechar{∪}{\ensuremath{\cup}}
\newunicodechar{⟨}{\ensuremath{\langle}}
\newunicodechar{⟩}{\ensuremath{\rangle}}
\newunicodechar{・}{\textperiodcentered}

\usepackage{xurl}
\usepackage{bookmark}
\usepackage{hyperref}

\setcounter{secnumdepth}{-\maxdimen}
\providecommand{\tightlist}{\setlength{\itemsep}{1pt}\setlength{\parskip}{0pt}}

\definecolor{openieAmber}{HTML}{D97706}
\definecolor{openieAmberBg}{HTML}{FFFBEB}
\definecolor{openieAmberBorder}{HTML}{F59E0B}
\definecolor{openieSlate}{HTML}{334155}
\definecolor{openieDark}{HTML}{0F172A}
\definecolor{openieCodeBg}{HTML}{F8FAFC}
\definecolor{openieCodeBorder}{HTML}{CBD5E1}
\definecolor{openieBlue}{HTML}{0057B8}

\hypersetup{
  pdftitle={RCE-VC: Law of Version Control},
  pdfauthor={David Charlot, Open Interface Engineering},
  colorlinks=true,
  linkcolor=openieBlue,
  citecolor=openieBlue,
  urlcolor=openieBlue,
  pdfcreator={XeTeX / Tectonic — Open Interface Engineering}
}

\urlstyle{same}
\setlength{\emergencystretch}{3em}
\hfuzz=6pt
\sloppy

\newcommand{\openiefig}[3]{%
  \IfFileExists{#1}{%
    \includegraphics[width=#2\linewidth,height=0.70\textheight,keepaspectratio]{#1}%
  }{%
    \fcolorbox{openieCodeBorder}{openieCodeBg}{%
      \parbox{0.92\linewidth}{\footnotesize\ttfamily Figure asset embedded in archival PDF; interactive/vector source: \url{#3}}%
    }%
  }%
}

\title{\vspace{-1.2em}\textbf{RCE-VC: Law of Version Control}}
\author{\textbf{David Charlot, Open Interface Engineering}\\[0.15em]
{\small Research study · draft $\cdot$ Open Interface Engineering $\cdot$ \href{https://research.openie.dev}{research.openie.dev}}}
\date{\small Revision \texttt{af57f48} (22 November 2026) $\cdot$ Web: \href{https://research.openie.dev/papers/rce-vc/}{research.openie.dev/papers/rce-vc/} $\cdot$ Archival LaTeX: \href{https://research.openie.dev/pdfs/latex/rce-vc.pdf}{/pdfs/latex/rce-vc.pdf} (\href{https://research.openie.dev/pdfs/latex/rce-vc.tex}{.tex})}

\begin{document}
\maketitle
\vspace{-0.6em}

% OpenIE Measurement Honesty Callout Box
\noindent\fcolorbox{openieAmberBorder}{openieAmberBg}{%
  \parbox{\dimexpr\linewidth-2\fboxsep-2\fboxrule\relax}{%
    \small\color{openieDark}%
    \textbf{Measurement.} Research study, draft, not a journal final. Milestones \texttt{M0} to \texttt{M2} are verified in public \texttt{openIE-dev/rce-vc} at commit \texttt{aba99fce} (\texttt{49} offline tests passing). Wall-clock times (\textbf{\texttt{69.7}} to \textbf{\texttt{73.6 \textmu s}}), admission-path shares, and Apple M5 Max per-process \texttt{L3} \texttt{reported\_j} (\textbf{\texttt{317.4}} to \textbf{\texttt{365.4 \textmu J}} per \texttt{100,000}-unit single-core admission) were measured on the reference build. Receipt costs also carry the stated \texttt{est\_uj} schedule. No hardware power rail or wall meter was attached (\texttt{measured\_j = None}, \texttt{board\_metered = false}). No FPGA board result. No fabricated citations.
  }%
}

\vspace{6pt}
% OpenIE Amber Left-Rule Summary Quote
\noindent{\color{openieAmber}\vrule width 3pt}\hspace{8pt}%
\parbox{\dimexpr\linewidth-11pt\relax}{%
  \small\itshape\color{openieSlate}%
  RCE-VC (a version-control reference, separate from the Reality Code Engine language) admits each change into one canonical repository state or refuses it with a recorded reason, splits Rust and Python files into name-free top-level items, and reads and writes ordinary Git history. Milestones M0, M1 and M2 are implemented at public commit aba99fce (49 offline tests). At 100,000 units on Apple M5 Max, admission takes 69.7 to 73.6 microseconds and 317 to 365 microjoules of per-process kernel rusage energy (reported\_j; measured\_j is unset).
}
\vspace{6pt}

{
\small
\setcounter{tocdepth}{2}
\tableofcontents
}
\vspace{4pt}
\hrule
\vspace{8pt}

\emph{Revision
\href{https://github.com/openIE-dev/research-openie-web/commits/main/src/content/papers/rce-vc.md}{\texttt{af57f48}},
22 November 2026. Reference implementation:
\href{https://github.com/openIE-dev/rce-vc}{openIE-dev/rce-vc} at
\href{https://github.com/openIE-dev/rce-vc/tree/aba99fcec4414a32bff30ab075668350d3fb80bf}{\texttt{aba99fcec441\allowbreak{}4a32bff30ab0\allowbreak{}75668350d3fb80bf}}
(13 November 2026), updating the initial \texttt{8bf9c9af} tree with
\texttt{rcevc-keys} top-level Rust/Python item splitting, the
\texttt{rcevc-\allowbreak{}core::assure} boundary (\texttt{49} offline tests), and
Apple M5 Max \texttt{reported\_j} admission measurements recorded on 22
November 2026.}

\section{Abstract}\label{abstract}

RCE-VC is a version-control design and a small reference in Rust. Many
writers, people or programs, submit changes against one shared state. A
named check admits the result, or the change is refused and the reason
is kept. Writers can stay on Git. The reference is not the Reality Code
Engine language. It uses the same framing: a Form, a total step, a
hash-chained log, name-free keys, and a ladder of gears with a model
last.

\textbf{In plain words}

\begin{enumerate}
\def\labelenumi{\arabic{enumi}.}
\tightlist
\item
  Any number of people and programs may write at once. A change enters
  the shared history when a named check accepts the result. If the check
  does not accept it, the change stays out and the reason is written
  down. Who wrote it does not decide.
\item
  Changes that do not touch the same things do not have to wait for one
  another. When they do overlap, a price written down in advance ranks
  the results that pass, and a result that was accepted once, including
  a resolved conflict, is reused. A model is asked only after the
  cheaper steps fail, and its answer still has to pass the check.
\item
  The reference reads and writes ordinary Git history and now splits
  Rust and Python files into name-free top-level items
  (\texttt{rcevc-keys}), so a writer need not leave Git. At
  \texttt{aba99fce} on Apple M5 Max, admitting a change against a
  100,000-unit repository took \texttt{69.7–73.6\ µs} and
  \texttt{317–365\ µJ} in operating-system per-process kernel
  \texttt{ri\_energy\_nj} (\texttt{reported\_j}). Receipt
  \texttt{est\_uj} fields remain a stated integer table
  (\texttt{424,153\ µJ} mean on the same workload), and no external
  hardware rail or wall monitor was attached
  (\texttt{measured\_\allowbreak{}j\ =\ None}, \texttt{board\_\allowbreak{}metered\ =\ false}).
\end{enumerate}

\section{Terms and labels}\label{terms-and-labels}

\begin{itemize}
\tightlist
\item
  SHA: SHA-256, the hash function the reference uses for keys, content
  digests and the log
\item
  BLAKE3: a cryptographic hash named as the design target for internal
  keys; the reference does not use it yet
\item
  HAMT: hash array mapped trie, the persistent map that holds repository
  state
\item
  UTF: UTF-8, the text encoding the Git file mapping accepts
\item
  MSRV: minimum supported Rust version (\texttt{1.85} for
  \texttt{rcevc-core} and \texttt{rcevc-proj-git}, \texttt{1.90} for
  \texttt{rcevc-keys})
\item
  CI: continuous integration
  (\texttt{cargo\ test\ -\allowbreak{}\/\allowbreak{}-\allowbreak{}locked\ -\allowbreak{}\/\allowbreak{}-\allowbreak{}offline\ -\allowbreak{}\/\allowbreak{}-\allowbreak{}all-\allowbreak{}targets} on
  \texttt{rustc\ 1.100.0})
\item
  CAS: compare-and-swap reference update (\texttt{CasRef} and
  \texttt{GitDir::compare\_\allowbreak{}and\_\allowbreak{}swap\_\allowbreak{}ref})
\item
  FIPS: Federal Information Processing Standards; the SHA-256 checks use
  the FIPS 180-4 vectors
\item
  RFC: Request for Comments; the inflate code follows RFC 1950 and RFC
  1951
\item
  SATB, CPU, GPU, PID, SMC, SoC: Satiation Benchmark telemetry
  discipline, central processing unit, graphics processing unit, process
  identifier, System Management Controller, and system-on-chip
\item
  VC1, VC2, VC3, VC4, VC5: the five propositions the reference tests,
  named with their tests below
\item
  CONGRA: a corpus named in the project notes for a future replay; the
  name is not expanded on this page
\item
  README: the overview file shipped in the reference repository
\end{itemize}

Three labels stay apart on this page across whole-chip (L2) and
per-process (L3) telemetry:

\begin{itemize}
\tightlist
\item
  \textbf{Measured (wall time and counts)} means a wall-clock time,
  instruction count or path count from the release build
  (\texttt{aba99fce} on Apple M5 Max, and the README baseline printed at
  \texttt{8bf9c9af}).
\item
  \textbf{Reported (\texttt{reported\_j}: L2 and L3 counters)} means an
  operating-system or on-chip counter reading on Apple M5 Max
  (\texttt{macOS\ 27.2.1}): whole-chip \texttt{IOReport} via
  \texttt{macmon} (L2, shared-host label \texttt{A2}, gated by Falsifier
  6 at \texttt{\textless{}=\ 5.0\%} pre/post idle drift) and per-PID
  kernel \texttt{proc\_\allowbreak{}pid\_\allowbreak{}rusage(\allowbreak{}RUSAGE\_\allowbreak{}INFO\_\allowbreak{}V6).\allowbreak{}ri\_\allowbreak{}energy\_\allowbreak{}nj}
  (L3, label \texttt{D}). No external hardware rail or wall-plug power
  probe was attached (\texttt{measured\_\allowbreak{}j\ =\ None},
  \texttt{board\_\allowbreak{}metered\ =\ false}).
\item
  \textbf{Stated (\texttt{est\_uj} / \texttt{est\_j})} means the default
  integer cost table copied onto receipts inside \texttt{rcevc-core}. It
  is an analytical table written in source code, never a counter or
  wall-meter reading.
\end{itemize}

\section{1. What the reference is, and what it is
not}\label{what-the-reference-is-and-what-it-is-not}

The project's aim is a history that many humans and agents can write at
once, without asking a writer to stop using Git. The framing in the
project notes is that branching, diffing and merging into one main line
is an old interface, that compatibility is what keeps that interface in
place, and that a deterministic rule should decide wherever a rule can,
with a model last.

The Reality Code Engine (RCE) is a programming language studied in its
own paper on this site. RCE-VC is not built on that language. It does
not parse RCE programs, emit them, or call the RCE checker. What it
shares is the framing: a Form for the repository, a total deterministic
step called δ, a hash-chained log whose fold is the state, name-free
keys, and a ladder of gears in which a model runs last and is not
trusted.

At commit \texttt{aba99fce} (\texttt{13\ November\ 2026}), the public
workspace is version 0.1.0 and holds \textbf{three} crates:

\begin{enumerate}
\def\labelenumi{\arabic{enumi}.}
\tightlist
\item
  \texttt{rcevc-core} is the deterministic core: no clock, no
  randomness, no I/O, no dependencies, and
  \texttt{\#!{[}forbid(unsafe\_code){]}} (MSRV \texttt{1.85}). It now
  also includes \texttt{rcevc-\allowbreak{}core::assure}, an assurance boundary
  ensuring that a checked action is minted only from an outcome already
  in that repository's log and applies only at the head and gate it was
  checked against, while \texttt{CasRef} retries one operation
  identifier so another writer's identical target is never mistaken for
  success.
\item
  \texttt{rcevc-keys} splits Rust (\texttt{.rs}) and Python
  (\texttt{.py}) source files into top-level items using vendored
  \texttt{tree-\allowbreak{}sitter\ 0.\allowbreak{}27.\allowbreak{}2}, \texttt{tree-\allowbreak{}sitter-\allowbreak{}rust\ 0.\allowbreak{}24.\allowbreak{}2}, and
  \texttt{tree-\allowbreak{}sitter-\allowbreak{}python\ 0.\allowbreak{}25.\allowbreak{}0} (MSRV \texttt{1.90}, built offline
  from \texttt{vendor/} with a C compiler and zero network fetches).
  Each item's identifier is the core name-free key of a normalized token
  body --- invariant to item renames, whitespace, and comments --- while
  exact original bytes stay in \texttt{parts} so \texttt{project}
  concatenates them back losslessly; non-\texttt{.rs}/\texttt{.py} paths
  or parse-error trees fall back cleanly to one leaf over the original
  bytes.
\item
  \texttt{rcevc-proj-git} is the Git projection
  (\texttt{\#!{[}forbid(unsafe\_code){]}}, MSRV \texttt{1.85}), carrying
  its own SHA-1, zlib and Git object code plus
  \texttt{GitDir::compare\_\allowbreak{}and\_\allowbreak{}swap\_\allowbreak{}ref} so the core stays free of
  dependencies.
\end{enumerate}

CI runs \texttt{cargo\ test\ -\allowbreak{}\/\allowbreak{}-\allowbreak{}locked\ -\allowbreak{}\/\allowbreak{}-\allowbreak{}offline\ -\allowbreak{}\/\allowbreak{}-\allowbreak{}all-\allowbreak{}targets}
and
\texttt{cargo\ clippy\ -\allowbreak{}\/\allowbreak{}-\allowbreak{}locked\ -\allowbreak{}\/\allowbreak{}-\allowbreak{}offline\ -\allowbreak{}\/\allowbreak{}-\allowbreak{}all-\allowbreak{}targets\ -\allowbreak{}\/\allowbreak{}-\allowbreak{}\ -\allowbreak{}D\ warnings}
on \texttt{rustc\ 1.100.0}, executing \textbf{49 offline tests}
(\texttt{29} in \texttt{rcevc-core}, \texttt{8} in \texttt{rcevc-keys},
and \texttt{12} in \texttt{rcevc-proj-git}, verified against
\texttt{git\ 2.43}). Propositions numbered past VC5 are not stated on
this page.

\section{2. Law of Version Control}\label{law-of-version-control}

The reference states it as follows. The same paragraph opens the
repository README and the crate documentation. It is quoted here once.

\begin{quote}
Version control is the admission of changes to a canonical, name-free
repository state. Each change is admitted only when a named checker
accepts the post-state under a declared contract, or is refused with a
recorded reason; changes whose dependency closures do not meet are
admitted without coordination; among admissible integrations, a stated
cost ranks the candidates; and every accepted result, including every
resolved conflict, is reused.
\end{quote}

The rest of this page is what the reference at \texttt{aba99fce}
implements, what was measured on it, and what it does not yet do. The
statement is the standard the reference is built to. It is not a report
that existing version-control tools already behave this way.

\section{3. Quantities}\label{quantities}

A \textbf{unit} is a body of text and a list of dependency slots. If the
dependency list is empty, the unit is a leaf and the body is literal
text, so a dollar sign in a shell script stays a dollar sign. If the
list is not empty, a dollar sign followed by an index selects that
dependency, and two dollar signs are a literal dollar sign.

The \textbf{key} of a unit is a name-free identity. A rename changes
which slot a name is bound to. It does not change the key. The reference
defines

\[\mathrm{key}(u) = H(\texttt{rcevc/unit/v1},\; \mathrm{body}(u),\; \mathrm{key}(d_0), \ldots, \mathrm{key}(d_{m-1}))\]

where \(H\) is SHA-256 through a versioned, domain-separated encoder,
and \(d_i\) are the dependency units. The \textbf{content digest}, which
the trie hashes, is a different value:
\(H(\mathrm{body}, \text{dependency names})\).

A \textbf{draft} is an unvalidated map, used for genesis and for
anything read from outside. Sealing a draft checks it and produces a
\textbf{state}. A state is well formed: dependencies resolve, every
index in a body is bound, and the dependency graph has no cycle. It
holds three persistent maps. The unit map is a Merkle HAMT. The key map
and the reverse-dependency index are persistent and are not part of the
Merkle root. Cloning a state and reading its root are \(O(1)\). Two
states are equal when their roots are equal. A diff walks only subtrees
whose hashes differ, so its cost follows the difference, not the size of
the repository. Dependents whose key changed, but whose own unit did
not, are left out of the diff on purpose. The closure, defined below,
carries those dependents. Putting them in the diff would turn a shared
dependent into a false conflict.

An \textbf{edit} changes slots and validates once at the end, over the
reverse-dependency closure of the slots touched. States inside an edit
may be ill formed. That is intentional, for example a deletion followed
by a put.

A \textbf{change} carries an opaque writer label, a base state, a level,
a list of operations, and a message. The operations are put, delete and
rename. The levels are name-only, erratum, meaning change and interface
change. The name-only level is checked by the kernel: the multiset of
keys must be unchanged. The change identifier covers every field. The
message is stated intent, for example a commit subject. Admission does
not read it, and it does not depend on the writer label.

For a change \(c\) against a state \(s\),

\[\mathrm{writes}(c, s) = \mathrm{diff}(s,\; \mathrm{apply}(s, c)).\]

Reads are the dependencies of puts and the sources of renames. The
closure of \(c\) in a state and in the post-state is the writes, the
dependents before, the dependents after, and the reads. Two changes
commute on \(s\) when each one's writes miss the other's closure:

\[\mathrm{commutes}(a, b, s) \iff \mathrm{writes}(a) \cap \mathrm{closure}(b) = \emptyset \;\land\; \mathrm{writes}(b) \cap \mathrm{closure}(a) = \emptyset.\]

The reference tests that commuting changes give the same root in either
order.

The step \textbf{δ} is \texttt{submit}. It is total on user input. It
returns an acceptance with a receipt, or a refusal with a reason, and it
appends one hash-chained log entry either way. The head is the fold of
that log. Replay verifies the chain and re-integrates every entry. It
uses the recorded resolution for a conflict. It does not call a
proposer. It checks each pre-state root, each post-state root, and each
refusal reason.

\section{4. The repository form}\label{the-repository-form}

The README lines the implemented form up as nine parts, plus the
\texttt{rcevc-keys} item splitter and \texttt{rcevc-\allowbreak{}core::assure}
boundary.

{\def\LTcaptype{none} % do not increment counter
{\footnotesize\setlength{\tabcolsep}{3.8pt}\renewcommand{\arraystretch}{1.15}
\begin{longtable}[]{@{}
  >{\raggedright\arraybackslash}p{(\linewidth - 4\tabcolsep) * \real{0.2083}}
  >{\raggedright\arraybackslash}p{(\linewidth - 4\tabcolsep) * \real{0.1979}}
  >{\raggedright\arraybackslash}p{(\linewidth - 4\tabcolsep) * \real{0.5938}}@{}}
\toprule\noalign{}
\begin{minipage}[b]{\linewidth}\raggedright
\textbf{Part}
\end{minipage} & \begin{minipage}[b]{\linewidth}\raggedright
\textbf{Where}
\end{minipage} & \begin{minipage}[b]{\linewidth}\raggedright
\textbf{What it does}
\end{minipage} \\
\midrule\noalign{}
\endhead
\bottomrule\noalign{}
\endlastfoot
R, name-free keys & \texttt{model.rs}, \texttt{rcevc-keys} & Keys as
above; \texttt{rcevc-keys} splits \texttt{.rs} and \texttt{.py} files
into top-level items with lossless byte projection. A rename changes no
key. \\
S, persistent state & \texttt{hamt.rs}, \texttt{model.rs} & Merkle HAMT.
Clone and root in \(O(1)\), diff in the size of the difference,
canonical trie shape. Keys and the reverse-dependency index update on
the edit, not on the whole repository. \\
Δ, typed changes & \texttt{change.rs} & Put, delete and rename against a
declared base, with a declared level. Name-only is checked by the
kernel. \\
X, commutation & \texttt{change.rs} & Writes are the actual diff. The
closure adds dependents before and after, and reads. Disjoint closures
commute. \\
K and G, checker and gate & \texttt{gate.rs} & Each checker states a
property class, a linking assumption, a trusted base, an authority, a
scope and a digest. \\
δ, the ladder & \texttt{repo.rs}, \texttt{assure.rs} & Lookup,
commutation, structural merge, re-validation, proposer, plus CAS ref
assurance (\texttt{CasRef}). The gate runs on every path. \\
D, the log & \texttt{repo.rs} & Hash-chained entries. The head is the
fold. Replay rebuilds it with no proposer. \\
U, reuse & \texttt{repo.rs} & A resolution is stored under a name-free
conflict signature, and only after the gate has accepted it. \\
C, cost & \texttt{repo.rs} & The stated per-rung table
(\texttt{est\_uj}). Receipts carry the sum. \\
\end{longtable}
}
}

A checker's authority is one of pinned, replay, proved, derived or
inferred. Its scope is one of three. \textbf{Unit}: re-check units in
the delta. \textbf{Closure}: the verdict depends on the unit and its
transitive dependencies, so the gate also re-checks dependents.
\textbf{Global}: check the whole state. If the previous head passed and
the declared scope is true, a scoped re-check matches a full check. The
test \texttt{delta\_\allowbreak{}gate\_\allowbreak{}agrees\_\allowbreak{}with\_\allowbreak{}full\_\allowbreak{}gate} records that
agreement over 600 trials, with more than 100 refusals. A checker that
lies about its scope breaks VC2. Scope is trusted. Nothing in this
version enforces it by hiding the rest of the state.

The gate digest covers every checker's statement, and that digest feeds
the genesis hash. The trie shape is canonical: a prefix holds nothing, a
single leaf, or a branch of two or more keys. Roots do not depend on the
order history arrived in. Everything hashed goes through the encoder
with a versioned domain tag of the form
\texttt{rcevc/\textless{}thing\textgreater{}/vN}. Changing what a tag
covers means bumping the version.

\section{5. The transition}\label{the-transition}

\texttt{submit} always takes one of the following paths. The gear names
are the ones written on the receipt.

\begin{enumerate}
\def\labelenumi{\arabic{enumi}.}
\tightlist
\item
  \textbf{G0, lookup.} The same change identifier against the same head,
  already refused, returns the stored refusal. The stated cost of that
  retry is the G0 row alone. The test
  \texttt{g0\_\allowbreak{}identical\_\allowbreak{}retry\_\allowbreak{}costs\_\allowbreak{}one\_\allowbreak{}lookup} records this.
\item
  \textbf{Base.} The base snapshot must exist, or the reason is an
  unknown base. Applying the operations to that base must succeed, or
  the reason is an apply error or an ill-formed unit. A name-only claim
  that does not leave the multiset of keys unchanged is refused.
\item
  \textbf{G1, commutation.} Let the moved set be the diff from the base
  to the current head. If the change's closure, in the base and in the
  head, misses that set, the change is applied to the head directly. No
  merge.
\item
  \textbf{G2, structural merge.} Otherwise the reference merges, slot by
  slot, over the union of what moved on the head and what the change
  itself changed. One side unchanged: take the other side. Both sides
  equal: take either. Any other slot is a conflict.
\item
  \textbf{Conflict.} The signature is
  \(H(\mathrm{key}_{base}, \mathrm{key}_{head}, \mathrm{key}_{theirs})\),
  which does not use the slot's name. G0 reuse returns a stored
  resolution for that signature. If none is stored, G4 asks the
  proposer. If there is no proposer, or it returns nothing, the change
  is refused with the conflicted slots. During replay the recorded
  resolution is used, and the proposer is not called.
\item
  \textbf{G3, re-validation.} The merged assignment is built on the head
  with a well-formedness check. If that check fails, the change is
  refused as ill formed.
\item
  \textbf{Gate, on every path, including G1.} The candidate and the diff
  from the current head are handed to the gate. A failing checker
  refuses the change and names the checker.
\item
  \textbf{Accept.} Proposer resolutions are copied into the reuse store
  only now. The new head is snapshotted. The receipt records the gears
  used, the gate evidence, the resolutions, and \texttt{est\_uj}.
\end{enumerate}

Refusals are logged too. The head does not move.

\section{6. Propositions the tests
hold}\label{propositions-the-tests-hold}

At \texttt{aba99fce}, \texttt{cargo\ test\ -\allowbreak{}\/\allowbreak{}-\allowbreak{}offline} runs 49 tests
(\texttt{29} core, \texttt{8} keys, \texttt{12} git), mapping the
propositions below:

{\def\LTcaptype{none} % do not increment counter
{\footnotesize\setlength{\tabcolsep}{3.8pt}\renewcommand{\arraystretch}{1.15}
\begin{longtable}[]{@{}
  >{\raggedright\arraybackslash}p{(\linewidth - 2\tabcolsep) * \real{0.4479}}
  >{\raggedright\arraybackslash}p{(\linewidth - 2\tabcolsep) * \real{0.5521}}@{}}
\toprule\noalign{}
\begin{minipage}[b]{\linewidth}\raggedright
\textbf{Claim}
\end{minipage} & \begin{minipage}[b]{\linewidth}\raggedright
\textbf{Test named in the repository}
\end{minipage} \\
\midrule\noalign{}
\endhead
\bottomrule\noalign{}
\endlastfoot
VC1, totality, no limbo: every submission is accepted or refused with a
reason, and logged &
\texttt{vc1\_\allowbreak{}every\_\allowbreak{}submission\_\allowbreak{}is\_\allowbreak{}accepted\_\allowbreak{}or\_\allowbreak{}refused\_\allowbreak{}and\_\allowbreak{}logged} \\
VC2, green by induction: every head is well formed and passes the gate &
\texttt{vc2\_\allowbreak{}every\_\allowbreak{}head\_\allowbreak{}is\_\allowbreak{}well\_\allowbreak{}formed\_\allowbreak{}and\_\allowbreak{}passes\_\allowbreak{}the\_\allowbreak{}gate} (400
random submissions, stale bases, poisoned changes) \\
VC3, commutation implies the same result in either order &
\texttt{vc3\_\allowbreak{}commutation\_\allowbreak{}implies\_\allowbreak{}order\_\allowbreak{}independence} (3,000 random
trials, more than 500 commuting pairs) and
\texttt{vc3\_\allowbreak{}disjoint\_\allowbreak{}writers\_\allowbreak{}integrate\_\allowbreak{}at\_\allowbreak{}g1\_\allowbreak{}in\_\allowbreak{}either\_\allowbreak{}order} \\
VC4, a semantic conflict is caught when the text does not conflict &
\texttt{vc4\_\allowbreak{}identical\_\allowbreak{}edits\_\allowbreak{}merge\_\allowbreak{}and\_\allowbreak{}semantic\_\allowbreak{}conflicts\_\allowbreak{}are\_\allowbreak{}caught\_\allowbreak{}at\_\allowbreak{}g3}
(an arity change against a new caller) \\
VC5, a resolved conflict is reused by its name-free signature &
\texttt{vc5\_\allowbreak{}resolved\_\allowbreak{}conflicts\_\allowbreak{}are\_\allowbreak{}reused\_\allowbreak{}by\_\allowbreak{}name\_\allowbreak{}free\_\allowbreak{}signature} \\
A model result is a candidate &
\texttt{model\_\allowbreak{}output\_\allowbreak{}is\_\allowbreak{}a\_\allowbreak{}candidate\_\allowbreak{}never\_\allowbreak{}trusted} \\
The state is the fold of the log &
\texttt{replay\_\allowbreak{}reproduces\_\allowbreak{}head\_\allowbreak{}without\_\allowbreak{}any\_\allowbreak{}proposer} \\
A rename changes no key, and name-only is checked &
\texttt{renames\_\allowbreak{}change\_\allowbreak{}no\_\allowbreak{}keys\_\allowbreak{}and\_\allowbreak{}name\_\allowbreak{}only\_\allowbreak{}is\_\allowbreak{}verified} \\
G0 idempotent retries cost one lookup &
\texttt{g0\_\allowbreak{}identical\_\allowbreak{}retry\_\allowbreak{}costs\_\allowbreak{}one\_\allowbreak{}lookup} \\
Scoped gate agrees with a full gate &
\texttt{delta\_\allowbreak{}gate\_\allowbreak{}agrees\_\allowbreak{}with\_\allowbreak{}full\_\allowbreak{}gate} (600 trials,
\textgreater100 refusals) \\
Snapshots untouched by later edits; incremental matches from-scratch &
\texttt{snapshots\_\allowbreak{}share\_\allowbreak{}structure\_\allowbreak{}and\_\allowbreak{}diff\_\allowbreak{}is\_\allowbreak{}exact}, HAMT tests
in \texttt{hamt.rs} \\
SHA-256 matches the published vectors & FIPS 180-4 vectors in
\texttt{hash.rs} \\
\end{longtable}
}
}

VC5's test accepts a conflict through the proposer, then submits the
same conflict under another slot name. The second acceptance takes the
G0 lookup gear, does not take G4, does not call the proposer again, and
has \texttt{est\_uj} less than one hundredth of the receipt from the
proposer path. That ratio is a comparison of stated receipt costs.

A separate test feeds the proposer a poisoned result. The gate refuses
it, the head is unchanged, and the refused resolution is not stored for
reuse.

\section{7. Measured wall time}\label{measured-wall-time}

The workload, as the swarm example implements it: every writer in a
round submits against that round's starting head, so later submissions
in the round are stale. One writer in five edits the hot module
(\texttt{below(5)\ ==\ 0}). One submission in fifty carries a forbidden
token (\texttt{below(\allowbreak{}50)\ ==\ 0}), which is the 2\% poison rate the
README states. The proposer in the example is a deterministic stand-in
that joins the two bodies. It is not a learned model. The example stores
10 units in each module. The fixed seed in the example is
\texttt{0x5eed}.

\subsection{7.1 Path shares at 400
units}\label{path-shares-at-400-units}

The README's path table is the default swarm: 64 writers, 20 rounds, 40
modules, 400 units, 1,280 submissions. Shares are of all submissions.

{\def\LTcaptype{none} % do not increment counter
{\footnotesize\setlength{\tabcolsep}{3.8pt}\renewcommand{\arraystretch}{1.15}
\begin{longtable}[]{@{}
  >{\raggedright\arraybackslash}p{(\linewidth - 2\tabcolsep) * \real{0.8229}}
  >{\raggedleft\arraybackslash}p{(\linewidth - 2\tabcolsep) * \real{0.1771}}@{}}
\toprule\noalign{}
\begin{minipage}[b]{\linewidth}\raggedright
\textbf{Path}
\end{minipage} & \begin{minipage}[b]{\linewidth}\raggedleft
\textbf{Share}
\end{minipage} \\
\midrule\noalign{}
\endhead
\bottomrule\noalign{}
\endlastfoot
Accepted at G1 (commuted, no merge) & 57.0\% \\
Accepted at G2, then G3 (structural merge, re-validated) & 27.3\% \\
Accepted at G2, then G4, then G3 (conflict, proposer, re-validated) &
13.8\% \\
Refused by the gate & 1.9\% \\
\end{longtable}
}
}

The four shares sum to 100. Every refusal was a poisoned change,
including 6 proposer outputs that inherited the poison. 1,259 of the
1,280 submissions were on a stale base, and 98.1\% of those stale
submissions were admitted with no writer rebase. The example made 182
proposer calls. Replay from genesis, with no proposer, reproduced the
same head.


\begin{center}
\openiefig{fig/rce-vc-fig01-swarm-and-timing.png}{0.94}{https://research.openie.dev/figures/rce-vc/swarm-and-timing.svg}
\end{center}


\emph{Figure 1. Measured path shares for the 400-unit, 20-round swarm,
and measured wall time per submission for the 5-round timing runs on the
baseline reference build. Release, one core. The stated receipt costs
are not on this figure.}

\subsection{7.2 Integration cost against repository
size}\label{integration-cost-against-repository-size}

The README's timing table is the same shape of workload with 64 writers
and 5 rounds, same seed, release build, one core. The cell is measured
wall time per submission. M0 is the earlier state that cloned maps. M1
is the persistent Merkle state.

{\def\LTcaptype{none} % do not increment counter
{\footnotesize\setlength{\tabcolsep}{3.8pt}\renewcommand{\arraystretch}{1.15}
\begin{longtable}[]{@{}
  >{\raggedleft\arraybackslash}p{(\linewidth - 8\tabcolsep) * \real{0.1237}}
  >{\raggedleft\arraybackslash}p{(\linewidth - 8\tabcolsep) * \real{0.1237}}
  >{\raggedleft\arraybackslash}p{(\linewidth - 8\tabcolsep) * \real{0.1856}}
  >{\raggedleft\arraybackslash}p{(\linewidth - 8\tabcolsep) * \real{0.2990}}
  >{\raggedleft\arraybackslash}p{(\linewidth - 8\tabcolsep) * \real{0.2680}}@{}}
\toprule\noalign{}
\begin{minipage}[b]{\linewidth}\raggedleft
\textbf{Units}
\end{minipage} & \begin{minipage}[b]{\linewidth}\raggedleft
\textbf{Modules}
\end{minipage} & \begin{minipage}[b]{\linewidth}\raggedleft
\textbf{M0, cloned maps}
\end{minipage} & \begin{minipage}[b]{\linewidth}\raggedleft
\textbf{M1, persistent state (README host)}
\end{minipage} & \begin{minipage}[b]{\linewidth}\raggedleft
\textbf{Ratio printed in the README}
\end{minipage} \\
\midrule\noalign{}
\endhead
\bottomrule\noalign{}
\endlastfoot
400 & 40 & 2,103 µs & 98 µs & 21 times \\
4,000 & 400 & 27,536 µs & 105 µs & 262 times \\
40,000 & 4,000 & not run & 136 µs & \\
100,000 & 10,000 & not run & 188 µs & \\
\end{longtable}
}
}

The 100,000-unit row is the example invoked as
\texttt{swarm\ -\allowbreak{}\/\allowbreak{}-\allowbreak{}\ 64\ 5\ 10000}: 64 writers, 5 rounds, 10,000
modules. M0 grew with repository size and M1 stayed near flat, with the
remaining rise coming from trie depth and from how much moved inside a
round, because snapshots share structure, diff and the gate skip equal
subtrees, an edit re-keys only edited units and their dependents, and
the gate re-checks the declared scope of the delta.

\subsection{\texorpdfstring{7.3 Measured wall time and on-chip
\texttt{reported\_j} at 100,000 units (commit
\texttt{aba99fce})}{7.3 Measured wall time and on-chip reported\_j at 100,000 units (commit aba99fce)}}\label{measured-wall-time-and-on-chip-reported_j-at-100000-units-commit-aba99fce}

On 22 November 2026, \texttt{openIE-\allowbreak{}dev/\allowbreak{}rce-\allowbreak{}vc} at commit
\texttt{aba99fce} (\texttt{rustc\ 1.100.0}, release build) was executed
and instrumented on Apple M5 Max (\texttt{macOS\ 27.2.1} \texttt{arm64})
under the dual-channel SATB \texttt{reported\_j} harness: whole-chip
\texttt{IOReport} via \texttt{macmon} (L2, label \texttt{A2}, Falsifier
6 \texttt{\textless{}=\ 5.0\%} idle-drift gate) paired with per-PID
kernel \texttt{proc\_\allowbreak{}pid\_\allowbreak{}rusage(\allowbreak{}RUSAGE\_\allowbreak{}INFO\_\allowbreak{}V6).\allowbreak{}ri\_\allowbreak{}energy\_\allowbreak{}nj} (L3,
label \texttt{D}), at \texttt{100,000} units (\texttt{10,000} modules ×
\texttt{10} units):

{\def\LTcaptype{none} % do not increment counter
{\scriptsize\setlength{\tabcolsep}{3.0pt}\renewcommand{\arraystretch}{1.14}
\begin{longtable}[]{@{}
  >{\raggedright\arraybackslash}p{(\linewidth - 12\tabcolsep) * \real{0.1933}}
  >{\raggedleft\arraybackslash}p{(\linewidth - 12\tabcolsep) * \real{0.1467}}
  >{\raggedleft\arraybackslash}p{(\linewidth - 12\tabcolsep) * \real{0.0600}}
  >{\raggedleft\arraybackslash}p{(\linewidth - 12\tabcolsep) * \real{0.1467}}
  >{\raggedleft\arraybackslash}p{(\linewidth - 12\tabcolsep) * \real{0.1533}}
  >{\raggedleft\arraybackslash}p{(\linewidth - 12\tabcolsep) * \real{0.1333}}
  >{\raggedleft\arraybackslash}p{(\linewidth - 12\tabcolsep) * \real{0.1667}}@{}}
\toprule\noalign{}
\begin{minipage}[b]{\linewidth}\raggedright
\textbf{\texttt{100,000}-unit workload (\texttt{aba99fce}, Apple M5
Max)}
\end{minipage} & \begin{minipage}[b]{\linewidth}\raggedleft
\textbf{Admissions / trial}
\end{minipage} & \begin{minipage}[b]{\linewidth}\raggedleft
\textbf{Trials}
\end{minipage} & \begin{minipage}[b]{\linewidth}\raggedleft
\textbf{Median wall time / admission}
\end{minipage} & \begin{minipage}[b]{\linewidth}\raggedleft
\textbf{Median L3 \texttt{reported\_j} / admission}
\end{minipage} & \begin{minipage}[b]{\linewidth}\raggedleft
\textbf{Instructions / admission}
\end{minipage} & \begin{minipage}[b]{\linewidth}\raggedleft
\textbf{Stated \texttt{CostTable} \texttt{est\_uj} / admission}
\end{minipage} \\
\midrule\noalign{}
\endhead
\bottomrule\noalign{}
\endlastfoot
\textbf{README swarm (\texttt{64} writers × \texttt{5} rounds)} &
\texttt{320} (\texttt{315} acc, \texttt{5} ref) & \texttt{5} &
\textbf{\texttt{69.69\ µs}} (\texttt{63.5–84.1\ µs}) &
\textbf{\texttt{365.35\ µJ}} (\texttt{346.4–395.7\ µJ}) &
\texttt{1,147,137} & \texttt{438,530.4\ µJ} \\
\textbf{Sustained window (\texttt{64} writers × \texttt{500} rounds)} &
\texttt{32,000} (\texttt{31,357} acc, \texttt{643} ref) & \texttt{3} &
\textbf{\texttt{73.56\ µs}} (\texttt{69.3–75.7\ µs}) &
\textbf{\texttt{317.21\ µJ}} (\texttt{278.9–369.0\ µJ}) &
\texttt{1,153,976} & \texttt{424,152.7\ µJ} \\
\end{longtable}
}
}

\begin{itemize}
\tightlist
\item
  \textbf{Genesis and path breakdown}: Building the
  \texttt{100,000}-unit Merkle HAMT genesis state took a median
  \textbf{\texttt{1.244\ s}} once (\texttt{5} trials). In the
  \texttt{32,000}-submission sustained run (\texttt{2.35\ s} window),
  \texttt{26,103} submissions (\texttt{81.6\%}) admitted at
  \texttt{G1Commute}, \texttt{2,594} (\texttt{8.1\%}) at
  \texttt{G2Structural\ \textgreater{}\ G3Revalidate}, \texttt{2,660}
  (\texttt{8.3\%}) via \texttt{G4Proposer} with \texttt{2,660} conflict
  signatures promoted for reuse, and \texttt{643} (\texttt{2.0\%})
  poisoned submissions were refused (\texttt{30,857} of \texttt{31,490}
  stale submissions = \texttt{98.0\%} admitted without writer rebase).
\item
  \textbf{L2 vs L3 consistency and Falsifier 6}: On the two sustained
  trials that cleared Falsifier 6 (\texttt{2.67\%} and \texttt{1.22\%}
  idle drift), whole-chip L2 \texttt{IOReport} net marginal energy was
  \texttt{9.22\ J} and \texttt{4.66\ J} (\texttt{288.2\ µJ} and
  \texttt{145.6\ µJ}/admission) beside per-PID L3
  \texttt{ri\_energy\_nj} of \texttt{11.81\ J} and \texttt{8.92\ J}
  (\texttt{369.0\ µJ} and \texttt{278.9\ µJ}/admission), while the
  middle trial (\texttt{7.28\%\ \textgreater{}\ 5.0\%} idle drift from
  concurrent host jobs) had its L2 net subtraction automatically
  invalidated by Falsifier 6. Notice that the placeholder
  \texttt{CostTable} (\texttt{424–439\ mJ}/admission) overstates
  CPU-only admission by roughly three orders of magnitude because its
  \texttt{5,000,000\ µJ} (\texttt{5\ J}) \texttt{Proposer} entry
  anticipates a neural proposer rather than the example's deterministic
  stand-in (\texttt{measured\_\allowbreak{}j\ =\ None},
  \texttt{board\_\allowbreak{}metered\ =\ false}).
\end{itemize}

\section{8. Stated cost on receipts}\label{stated-cost-on-receipts}

The default table in \texttt{CostTable} at \texttt{aba99fce}, in
microjoules, remains a stated price per row inside repository receipts:

{\def\LTcaptype{none} % do not increment counter
{\footnotesize\setlength{\tabcolsep}{3.8pt}\renewcommand{\arraystretch}{1.15}
\begin{longtable}[]{@{}
  >{\raggedright\arraybackslash}p{(\linewidth - 4\tabcolsep) * \real{0.1667}}
  >{\raggedright\arraybackslash}p{(\linewidth - 4\tabcolsep) * \real{0.6458}}
  >{\raggedleft\arraybackslash}p{(\linewidth - 4\tabcolsep) * \real{0.1875}}@{}}
\toprule\noalign{}
\begin{minipage}[b]{\linewidth}\raggedright
\textbf{Row}
\end{minipage} & \begin{minipage}[b]{\linewidth}\raggedright
\textbf{What adds it}
\end{minipage} & \begin{minipage}[b]{\linewidth}\raggedleft
\textbf{Stated µJ}
\end{minipage} \\
\midrule\noalign{}
\endhead
\bottomrule\noalign{}
\endlastfoot
G0 & Charged at the start of every attempt. Charged again for each
reused conflict signature. A memoised refusal is this row alone. & 1 \\
G1 & Charged when the attempt reaches the commutation test & 10 \\
G2 & Charged when the attempt performs the structural merge & 100 \\
Checker & Charged once per checker, on every path that reaches the gate,
including G1 & 1,000 \\
Proposer & Charged once per proposer call & 5,000,000 \\
\end{longtable}
}
}

G3 is written on the path and has no row of its own. A receipt's
\texttt{est\_uj} is the sum of the rows that attempt added. Replacing
this table with hardware-rail \texttt{measured\_j} is milestone M6 and
remains open; Section 7.3 provides the first operating-system
\texttt{reported\_j} calibration (\texttt{317–365\ µJ} per
\texttt{100,000}-unit admission with a deterministic stand-in proposer).

\section{9. The Git projection}\label{the-git-projection}

The projection is milestone M2, in \texttt{rcevc-proj-git}. A leaf unit
is the file bytes. A unit with dependencies is written as a file with a
short header, which the reference treats as a bijection on UTF-8 files:

\begin{verbatim}
#rcevc-unit v1
#dep <slot>
#body
<body>
\end{verbatim}

One \texttt{\#dep} line is written for each dependency.

\textbf{Out.} \texttt{export} writes one commit on \texttt{main} for
each accepted entry. The author is the writer label. The subject is the
change message, or a line naming the change identifier and the gears if
the message is empty. Each accepted commit carries the trailers
\texttt{RCE-VC-Change}, \texttt{RCE-VC-Pre}, \texttt{RCE-VC-Post},
\texttt{RCE-VC-Path}, \texttt{RCE-VC-Evidence},
\texttt{RCE-VC-Resolved}, \texttt{RCE-VC-Est-uJ} and
\texttt{RCE-VC-Entry}. \texttt{RCE-VC-Est-uJ} is the stated microjoule
sum, not a measurement. The ref \texttt{refs/rcevc/log} has one commit
per log entry, refusals included, and a refusal commit also carries
\texttt{RCE-VC-Reason}. Given the same log and the same time base,
export produces the same commit identifiers.

\textbf{In.} \texttt{ingest} walks parent commits from a tip until it
finds a commit whose \texttt{RCE-VC-Post} trailer is a snapshot. That
snapshot is the base. The tip tree becomes one change, put and delete,
against that base. If the content multiset is unchanged, the level is
name-only. The branch is the proposal. The ladder decides. The writer is
not asked to rebase.

The README records these checks against git 2.43, in
\texttt{tests/\allowbreak{}git\_\allowbreak{}interop.\allowbreak{}rs} (\texttt{12} tests in
\texttt{rcevc-proj-git} at \texttt{aba99fce}):

{\def\LTcaptype{none} % do not increment counter
{\footnotesize\setlength{\tabcolsep}{3.8pt}\renewcommand{\arraystretch}{1.15}
\begin{longtable}[]{@{}
  >{\raggedright\arraybackslash}p{(\linewidth - 2\tabcolsep) * \real{0.5208}}
  >{\raggedright\arraybackslash}p{(\linewidth - 2\tabcolsep) * \real{0.4792}}@{}}
\toprule\noalign{}
\begin{minipage}[b]{\linewidth}\raggedright
\textbf{Check}
\end{minipage} & \begin{minipage}[b]{\linewidth}\raggedright
\textbf{Result recorded in the README}
\end{minipage} \\
\midrule\noalign{}
\endhead
\bottomrule\noalign{}
\endlastfoot
\texttt{git\ fsck\ -\allowbreak{}\/\allowbreak{}-\allowbreak{}full\ -\allowbreak{}\/\allowbreak{}-\allowbreak{}strict} on the exported history &
clean \\
Every commit on \texttt{main} re-imports to the root in its
\texttt{RCE-VC-Post} trailer & yes \\
Two exports of the same log in different directories & identical commit
identifiers \\
Git's trailer parser reads the receipts & yes \\
Two plain Git commits on a branch whose base has gone stale & ingested
as one change, admitted at G1, \texttt{main} moves forward \\
Objects Git itself wrote, including dynamic-Huffman deflate & read \\
A plain Git commit that contains a forbidden token & refused, recorded
on \texttt{refs/rcevc/log}, \texttt{main} unmoved \\
\end{longtable}
}
}

Limits of this projection, as the README states them: loose objects and
loose refs only; file mode 100644 on the way out, so the executable bit
is not yet a unit attribute; symlinks, submodules and non-UTF-8 files
are refused with a reason rather than guessed.

The project notes name Mergiraf, a structured merge that uses
tree-sitter, as work aimed at parse-shaped files. That is a convergence
observation. This page does not score the reference against it, and it
does not score the reference against any other version-control system.

\section{10. Limits of this version}\label{limits-of-this-version}

At commit \texttt{aba99fce}:

\begin{enumerate}
\def\labelenumi{\arabic{enumi}.}
\tightlist
\item
  \texttt{rcevc-keys} splits Rust (\texttt{.rs}) and Python
  (\texttt{.py}) files into top-level items (with dependencies pointing
  only at earlier items in the same file), but TypeScript splitting and
  a 50-package / 1,000-repository round-trip evaluation are not yet
  done; other file types remain whole-file leaf units.
\item
  Git storage is loose objects and loose refs. There are no packfiles
  and no packed refs. Export writes mode 100644 only. Symlinks,
  submodules and non-UTF-8 files are refused.
\item
  Each exported commit rebuilds the whole tree.
\item
  The reverse-dependency sets are cloned when they are updated. A unit
  with many dependents makes that update large.
\item
  Key computation is recursive. A dependency chain of more than about
  10,000 units could overflow the stack.
\item
  One full-content root is linear in the number of units. Ingest uses it
  only to guess the name-only level. The kernel's name-only check walks
  the diff.
\item
  Snapshots are kept in memory, with no retention policy and no separate
  store.
\item
  \texttt{submit} is single-threaded (\texttt{CasRef} guards ref swaps,
  while \texttt{submit} admits commuting changes sequentially). Leases
  are not built.
\item
  A checker is trusted to tell the truth about its scope (Wasmtime
  sandboxing M3 is not yet built).
\item
  On-chip \texttt{reported\_j} (\texttt{317–365\ µJ} per
  \texttt{100,000}-unit admission) has been recorded externally via
  kernel \texttt{ri\_energy\_nj}, while in-repo receipt fields still
  carry the stated \texttt{est\_uj} table and external hardware-rail
  \texttt{measured\_j} (M6) remains open (\texttt{measured\_\allowbreak{}j\ =\ None},
  \texttt{board\_\allowbreak{}metered\ =\ false}).
\end{enumerate}

\section{11. Next work}\label{next-work}

\begin{itemize}
\tightlist
\item
  \textbf{Complete M1 (\texttt{rcevc-keys} expansion and corpus round
  trip).} Add the TypeScript grammar alongside Rust and Python, support
  cross-file item dependencies, and verify byte-identical round trips
  across at least 50 real crates and packages with independent edits to
  separate functions of one file commuting at G1.
\item
  \textbf{Complete M2 (BLAKE3 internal keys and \texttt{gix}
  packfiles).} Move internal hashing to BLAKE3 (keeping FIPS 180-4
  SHA-256 vectors as their own test and bumping domain tags) and move
  Git object storage to \texttt{gix} (\texttt{gitoxide}) for packfiles
  and packed refs, followed by a 1,000-repository history round trip.
\item
  \textbf{M3 (\texttt{rcevc-gate} on Wasmtime).} Run checkers as
  fuel-limited, deterministic WebAssembly modules with pinned digests so
  declared scope is enforced by exposing only the declared slice.
\item
  \textbf{M4 (leases and monotone channel).} Concurrent streaming for
  monotone changes without coordination.
\item
  \textbf{M5 (conflict corpora replay).} Replay Merge-Bench, CONGRA and
  AgenticFlict through the ladder, counting resolutions at G0, G1, G2,
  G3 and G4.
\item
  \textbf{M6 (\texttt{rcevc-meter}).} Replace \texttt{est\_uj} with
  calibrated \texttt{reported\_j} and external hardware-rail
  \texttt{measured\_j}.
\end{itemize}

\section{12. Copyright and license}\label{copyright-and-license}

Copyright 2026 Open Interface Engineering, Inc.

The reference implementation in
\href{https://github.com/openIE-dev/rce-vc}{openIE-dev/rce-vc} at commit
\texttt{aba99fcec441\allowbreak{}4a32bff30ab0\allowbreak{}75668350d3fb80bf}
(\texttt{13\ November\ 2026}) is under the Apache License, Version 2.0.
The license file in that tree states the copyright as Open Interface
Engineering, Inc., and the crate manifest sets
\texttt{license\ =\ "Apache-\allowbreak{}2.\allowbreak{}0"}. This page is the study text for that
reference.


\end{document}
